% ============================================================
%  第二章 曲面的第一基本形式
%  结构沿袭苏步青、胡和生《微分几何》第二章
%  记号约定：位置向量/切向量用 \boldsymbol{}，单位法向量用
%            \mathbf{n}，标量系数（E,F,G 等）不加粗。
% ============================================================

\section{曲面的概念}

\begin{definition}[正则曲面]
	设 $U$ 是 $\mathbb{R}^2$ 中的一个开区域，映射
	\begin{equation}
		\boldsymbol{r}:U\longrightarrow\mathbb{E}^3,\qquad
		(u,v)\longmapsto\boldsymbol{r}(u,v)
		\label{eq:曲面:定义}
	\end{equation}
	是 $C^k$（$k\geqslant2$）的，且在其上处处满足
	\begin{equation}
		\boldsymbol{r}_u(u,v)\times\boldsymbol{r}_v(u,v)\neq\boldsymbol{0},
		\label{eq:曲面:正则条件}
	\end{equation}
	其中 $\boldsymbol{r}_u=\partial\boldsymbol{r}/\partial u$、$\boldsymbol{r}_v=\partial\boldsymbol{r}/\partial v$，则称 $\boldsymbol{r}(U)$ 为\textbf{正则参数曲面}，$u$、$v$ 称为曲面的参数。
\end{definition}

正则条件式~\eqref{eq:曲面:正则条件} 等价于两个偏导向量 $\boldsymbol{r}_u,\boldsymbol{r}_v$ 处处线性无关，即矩阵 $(\boldsymbol{r}_u,\boldsymbol{r}_v)$ 的秩处处为 $2$。

\begin{definition}[切平面与法向量]
	在正则点处，由 $\boldsymbol{r}_u$、$\boldsymbol{r}_v$ 张成的平面称为曲面在该点的\textbf{切平面}；单位法向量定义为
	\begin{equation}
		\mathbf{n}=\frac{\boldsymbol{r}_u\times\boldsymbol{r}_v}
		{\left|\boldsymbol{r}_u\times\boldsymbol{r}_v\right|}.
		\label{eq:曲面:单位法向量}
	\end{equation}
\end{definition}

\begin{example}[球面]
	单位球面可由参数表示
	\begin{equation}
		\boldsymbol{r}(u,v)=\bigl(\cos u\cos v,\;\cos u\sin v,\;\sin u\bigr),
		\qquad -\frac{\pi}{2}\leqslant u\leqslant\frac{\pi}{2},\quad 0\leqslant v<2\pi,
		\label{eq:球面:参数表示}
	\end{equation}
	其中 $u$ 为纬度，$v$ 为经度。可以验证 $\boldsymbol{r}_u\times\boldsymbol{r}_v=\cos u\,\boldsymbol{r}\neq\boldsymbol{0}$（$u\neq\pm\pi/2$），故除两极外是正则曲面。
\end{example}

\section{参数曲线网}

\begin{definition}[坐标曲线]
	固定一个参数而让另一个参数变化所得的曲线称为\textbf{坐标曲线}：
	\begin{equation}
		u\text{曲线（$v$ 曲线）：}\boldsymbol{r}=\boldsymbol{r}(u,v_0),\qquad
		v\text{曲线（$u$ 曲线）：}\boldsymbol{r}=\boldsymbol{r}(u_0,v),
		\label{eq:坐标曲线:定义}
	\end{equation}
	全体坐标曲线构成曲面上的\textbf{参数曲线网}。
\end{definition}

\begin{definition}[参数变换]
	设 $(\bar u,\bar v)\mapsto(u,v)$ 是开区域之间的 $C^k$ 双射，其 Jacobi 行列式处处不为零：
	\begin{equation}
		\frac{\partial(u,v)}{\partial(\bar u,\bar v)}
		=
		\begin{vmatrix}
			u_{\bar u} & u_{\bar v}\\
			v_{\bar u} & v_{\bar v}
		\end{vmatrix}
		\neq0,
		\label{eq:参数变换:Jacobi}
	\end{equation}
	则称 $(\bar u,\bar v)$ 是曲面的一组新的参数，这样的变换称为\textbf{参数变换}（正则变换）。在参数变换下，曲面与切平面、法向量均不变。
\end{definition}

\section{第一基本形式}

\begin{definition}[第一基本形式]
	曲面的第一基本形式（度量形式）定义为切向量自身的点积：
	\begin{equation}
		\mathrm{I}=\mathrm{d}\boldsymbol{r}\cdot\mathrm{d}\boldsymbol{r}
		=E\,\mathrm{d}u^{2}+2F\,\mathrm{d}u\,\mathrm{d}v+G\,\mathrm{d}v^{2},
		\label{eq:第一基本形式:定义}
	\end{equation}
	其中系数
	\begin{equation}
		E=\boldsymbol{r}_u\cdot\boldsymbol{r}_u,\qquad
		F=\boldsymbol{r}_u\cdot\boldsymbol{r}_v,\qquad
		G=\boldsymbol{r}_v\cdot\boldsymbol{r}_v.
		\label{eq:第一基本形式:系数}
	\end{equation}
\end{definition}

\begin{theorem}[第一基本形式的正定性]
	第一基本形式是正定的二次微分形式：
	\begin{equation}
		EG-F^{2}=\left|\boldsymbol{r}_u\times\boldsymbol{r}_v\right|^{2}>0,\qquad E>0,\;G>0.
		\label{eq:第一基本形式:正定性}
	\end{equation}
\end{theorem}

\begin{remark}
	第一基本形式刻画了曲面上的度量（内蕴几何），它不随曲面的空间位置而变，只依赖于曲面上弧长的度量。这是"曲面论是二维黎曼几何"的出发点。
\end{remark}

\section{曲面上曲线的弧长}

\begin{proposition}[弧长公式]
	设曲面上的曲线为 $u=u(t)$，$v=v(t)$（$a\leqslant t\leqslant b$），则它的弧长为
	\begin{equation}
		L=\int_{a}^{b}\sqrt{E\dot{u}^{2}+2F\dot{u}\dot{v}+G\dot{v}^{2}}\;\mathrm{d}t,
		\label{eq:曲面弧长:公式}
	\end{equation}
	其中 $\dot u=\mathrm{d}u/\mathrm{d}t$、$\dot v=\mathrm{d}v/\mathrm{d}t$。弧长元素为
	\begin{equation}
		\mathrm{d}s^{2}=E\,\mathrm{d}u^{2}+2F\,\mathrm{d}u\,\mathrm{d}v+G\,\mathrm{d}v^{2}.
		\label{eq:曲面弧长:弧长元素}
	\end{equation}
\end{proposition}

\section{两方向间的夹角与正交参数曲线网}

\begin{definition}[两方向间的夹角]
	设曲面上一点的两个方向分别由 $(\mathrm{d}u,\mathrm{d}v)$ 与 $(\delta u,\delta v)$ 给出，则它们之间的夹角 $\theta$ 由下式决定：
	\begin{equation}
		\cos\theta=
		\frac{E\,\mathrm{d}u\,\delta u+F(\mathrm{d}u\,\delta v+\mathrm{d}v\,\delta u)+G\,\mathrm{d}v\,\delta v}
		{\sqrt{E\,\mathrm{d}u^{2}+2F\,\mathrm{d}u\,\mathrm{d}v+G\,\mathrm{d}v^{2}}
		 \sqrt{E\,\delta u^{2}+2F\,\delta u\,\delta v+G\,\delta v^{2}}}.
		\label{eq:两方向夹角:公式}
	\end{equation}
\end{definition}

\begin{proposition}[正交参数曲线网]
	$u$ 曲线与 $v$ 曲线正交当且仅当 $F=0$。事实上，令 $(\mathrm{d}u,\mathrm{d}v)=(1,0)$、$(\delta u,\delta v)=(0,1)$ 代入式~\eqref{eq:两方向夹角:公式}，得
	\begin{equation}
		\cos\theta=\frac{F}{\sqrt{EG}},
		\label{eq:正交曲线网:条件}
	\end{equation}
	故 $u$ 曲线与 $v$ 曲线正交 $\iff F=0$。$F=0$ 的参数曲线网称为\textbf{正交参数曲线网}。
\end{proposition}

\begin{remark}
	可以证明：在正则曲面上总能选取正交参数曲线网（例如借助等温参数）。本章及下一章中，许多公式在正交参数下会显著简化。
\end{remark}

\section{等距变换与保形变换}

\begin{definition}[等距变换]
	曲面 $S_1$ 与 $S_2$ 之间的一个 $C^1$ 一一对应称为\textbf{等距变换}（保长对应），如果它保持曲面上任意曲线的长度不变。
\end{definition}

\begin{theorem}[等距变换的判别]
	两个曲面是等距的，当且仅当可以选取适当的参数，使两者的第一基本形式完全相同，即
	\begin{equation}
		E\,\mathrm{d}u^{2}+2F\,\mathrm{d}u\,\mathrm{d}v+G\,\mathrm{d}v^{2}
		=\bar E\,\mathrm{d}\bar u^{2}+2\bar F\,\mathrm{d}\bar u\,\mathrm{d}\bar v+\bar G\,\mathrm{d}\bar v^{2}.
		\label{eq:等距变换:判别}
	\end{equation}
\end{theorem}

\begin{example}
	圆柱面与平面是等距的（把圆柱面沿母线剪开摊平即可），但球面与平面在局部上都不存在等距变换——这一点由高斯曲率的性质（见第三章）可以严格说明。
\end{example}

\begin{definition}[保形变换]
	曲面之间的一个一一对应称为\textbf{保形变换}（保角变换、共形变换），如果它保持曲面上任意两条曲线的交角不变。保形变换将任意一个小圆变为小圆。
\end{definition}

\begin{theorem}[保形变换的判别]
	两个曲面之间存在保形变换，当且仅当可以选取适当的参数使两者的第一基本形式成比例：
	\begin{equation}
		E:F:G=\bar E:\bar F:\bar G.
		\label{eq:保形变换:判别}
	\end{equation}
\end{theorem}

\begin{remark}
	由式~\eqref{eq:正交曲线网:条件} 可见夹角公式只用到 $E,F,G$ 的比值，故第一基本形式相差一个正函数因子的曲面之间必存在保形变换。平面到单位球面（例如经 Mercator 投影）即是一例。
\end{remark}

\section{曲面上的面积}

\begin{theorem}[面积公式]
	设 $D$ 是参数区域，则曲面上对应于 $D$ 的面积为
	\begin{equation}
		A=\iint_{D}\left|\boldsymbol{r}_u\times\boldsymbol{r}_v\right|\,\mathrm{d}u\,\mathrm{d}v
		=\iint_{D}\sqrt{EG-F^{2}}\;\mathrm{d}u\,\mathrm{d}v.
		\label{eq:面积:公式}
	\end{equation}
\end{theorem}

\begin{proof}
	微元平行四边形由 $\boldsymbol{r}_u\,\mathrm{d}u$ 与 $\boldsymbol{r}_v\,\mathrm{d}v$ 张成，其面积为
	\begin{equation}
		\left|\boldsymbol{r}_u\,\mathrm{d}u\times\boldsymbol{r}_v\,\mathrm{d}v\right|
		=\left|\boldsymbol{r}_u\times\boldsymbol{r}_v\right|\,\mathrm{d}u\,\mathrm{d}v,
		\label{eq:面积:微元}
	\end{equation}
	再利用式~\eqref{eq:第一基本形式:正定性} 即得。面积微元 $\mathrm{d}A=\sqrt{EG-F^{2}}\,\mathrm{d}u\,\mathrm{d}v$ 是内蕴量。
\end{proof}
